
\chapter{ISO 8583: язык, на котором говорят платежи}\label{ch:iso8583}
В~процессинге лежит авторизационный запрос на покупку на сумму
1\,250 рублей 50 копеек:

\begin{center}
  \small
  \texttt{30 31 30 30}\\[2pt]
  \texttt{72 3C 04 81 20 C0 80 00}\\[2pt]
  \texttt{31 36 32 32 30 30 31 32 31 32 33 34 35 36 30 38 39 32}\\[2pt]
  \texttt{\ldots}
\end{center}

Чтобы превратить дамп в~рассказ о~транзакции, нужно прочитать тип сообщения, затем карту присутствующих полей, затем значения по порядку.
Ниже разобран MTI~0100 и~его ответ MTI~0110~-- тот же жизненный цикл, что в~главе~\ref{ch:tx-lifecycle}, только в~байтах.

\section{Одно сообщение в~логах}\label{sec:iso8583-structure}

Грамматика ISO~8583 одинакова для всех сообщений и состоит из трёх частей, идущих строго в~одном порядке:

\begin{enumerate}
  \item \textbf{MTI} (Message Type Indicator): четыре цифры,
    которые отвечают на вопрос, \emph{что это за сообщение}:
    авторизация, финансовое сообщение, сообщение отмены,
    сетевое
    управление.
  \item \textbf{Bitmap}: битовая карта, которая отвечает
    на вопрос, \emph{какие поля присутствуют дальше}.
  \item \textbf{Элементы данных} (data elements, DE): сами значения: номер карты
    (PAN, Primary Account Number; разбор~--- в~гл.~\ref{ch:tx-lifecycle}),
    сумма, STAN (Systems Trace Audit Number, номер трассировки),
    код ответа, идентификатор терминала
    и десятки других признаков транзакции.
\end{enumerate}

\highlight{ISO~8583 устроен иначе, чем самоописывающиеся форматы вроде JSON или XML.
Сначала идёт тип операции, затем bitmap с~картой присутствия, затем значения подряд в~порядке возрастания номеров}. Тег и~длина внутри каждого поля не передаются.

Канонический пример~-- MTI~0100, запрос на авторизацию, и~его пара MTI~0110, ответ эмитента.
Байты в~запросе и~ответе различаются, а~грамматика общая, так что навык чтения переносится с~одного MTI на другой.

\section{MTI -- где сообщение находится в~жизненном цикле}\label{sec:iso8583-mti}

MTI~-- первые четыре цифры сообщения. Они определяют тип операции: запрос на одобрение, финансовую запись, отмену, служебную проверку связи (echo-test).

Учебные обзоры раскладывают MTI по разрядам, но на практике надёжнее смотреть в~список разрешённых значений из спецификации интерфейса партнёра (IFS, Interface Specification; подробно~--- в~§\,\ref{sec:iso8583-implementations}).
Стандарт задаёт общую схему, конкретный процессинг или сеть фиксирует собственный набор MTI и~правил их использования; даже публичные профили расходятся в~деталях~\cite{iso-8583-2023,fis-iso8583-guide-2023}.

Авторизационный и~служебный потоки используют следующие MTI.

\medskip
\noindent
\begingroup
\small
\setlength{\tabcolsep}{4pt}
\renewcommand{\arraystretch}{1.2}
\begin{tabularx}{\textwidth}{
    @{}P{0.12\textwidth}
    Y@{}
  }
  \toprule
  \headrow
  \thd{MTI} & \thd{Смысл} \\
  \midrule
  MTI~0100 & Авторизационный запрос \\
  MTI~0110 & Ответ на авторизацию \\
  MTI~0120 & Authorization Advice -- уведомление об авторизации,
  принятой офлайн или через stand-in; гарантированная доставка
  без ожидания ответа \\
  MTI~0121 & Повтор Authorization Advice при недоставке MTI~0120 \\
  MTI~0200 & Финансовый запрос \\
  MTI~0210 & Ответ на финансовый запрос \\
  MTI~0400 & Acquirer Reversal Request -- запрос отмены от эквайера \\
  MTI~0410 & Ответ на отмену (reversal) \\
  MTI~0420 & Acquirer Reversal Advice -- одностороннее уведомление об отмене \\
  MTI~0430 & Ответ на reversal advice \\
  MTI~0800 & Сетевой служебный запрос: sign-on, echo, обмен ключами \\
  MTI~0810 & Ответ на сетевой служебный запрос \\
  \bottomrule
\end{tabularx}%
\endgroup
\medskip
Пара MTI~0100/MTI~0110 описывает двухстадийную авторизацию, когда эмитент принял решение, а~деньги ещё не двинулись в~расчётном смысле.
MTI~0120~-- advice: ответа на него не предполагается, доставка гарантирована. Типичное применение~-- deferred auth и уведомления stand-in processing.
Пара MTI~0200/MTI~0210~-- финансовое сообщение, и~во многих карточных диалектах именно она реализует одностадийный сценарий.
Отмена идёт через MTI~0400/MTI~0410 в~одних интерфейсах и через MTI~0420/MTI~0430 в~других. Какая пара применяется, определяется спецификацией интерфейса~\cite{fis-iso8583-guide-2023}.

Служебные пары MTI~0800/MTI~0810 обслуживают \emph{сетевой диалог}: \enquote{sign-on}~--- логический вход в~сеть с~началом сессии и~обновлением ключей; \enquote{echo} (\enquote{echo-test})~--- периодическая проверка живости канала; \enquote{cutover}~--- переключение на новый набор ключей или новую версию профиля без разрыва связи.

\subsection*{MTI~0100 -- запрос на решение}

MTI~0100~-- запрос с~данными, достаточными, чтобы эмитент решил, одобрять ли операцию: кто платит, сколько, где и~каким способом считана карта.
Минимальный набор включает PAN или Track~2, код обработки (processing code), сумму, дату и~время передачи, STAN, режим ввода карты, идентификаторы терминала и ТСП. Для чиповой операции добавляется DE~55~\cite{fis-iso8583-guide-2023}.

\subsection*{MTI~0110 -- ответ эмитента}

MTI~0110~-- решение эмитента. Ответ короче запроса, потому что PAN и Track~2 не нужны, зато появляются DE~39 (response code) и, при одобрении, DE~38 (authorization ID response).
STAN, сумма, дата и идентификатор терминала повторяются ради сопоставления ответа с~запросом~\cite{fis-iso8583-guide-2023}.

\subsection*{Отмена (reversal) -- MTI~0400 и MTI~0420}

Одни интерфейсы оформляют отмену как MTI~0400/MTI~0410, другие как MTI~0420/MTI~0430. FIS/Worldpay разделяет варианты request и~advice~\cite{fis-iso8583-guide-2023}.
Смысл общий~-- \enquote{считать предыдущее сообщение несостоявшимся} или \enquote{скорректировать его последствия}. Возврат денег клиенту~-- отдельная операция.

Сообщение об отмене должно однозначно сослаться на исходную операцию. Для этого используются тот же STAN, тот же временной контекст, а~в~ряде профилей~-- специальные \enquote{original data elements} (поля, в~которые упаковываются ссылки на исходную операцию: её MTI, STAN, дата и~время) вроде DE~90. Официальные профили прямо описывают DE~90 как способ передать эмитенту значения исходной операции~\cite{fis-iso8583-guide-2023}.

\practicenote{%
  Одинаковый MTI~0100 в~двух разных IFS даёт разные запросы. Один
  партнёр считает DE~32 обязательным, другой не принимает его вовсе;
  у~одного DE~22 кодирует бесконтактный сценарий кодом \texttt{071},
  у~другого~--- \texttt{05F}. Код, работающий с~одним процессингом,
  молча падает на неожиданном поле у~другого. Диалектные различия
  ловить надо на тестовых транзакциях; в~production они обходятся дороже.
}

\section{Bitmap -- какие поля искать дальше}\label{sec:iso8583-bitmap}

\highlight{После MTI идёт bitmap~-- восемь или шестнадцать байт, в~которых каждый бит отвечает за присутствие одного элемента данных}.
Сами значения bitmap не несёт; он показывает только, какие поля искать дальше в~теле сообщения~\cite{fis-iso8583-guide-2023}.

Bitmap~--- лишь одна из возможных структур. Поздние форматы вроде TLV (Tag-Length-Value), как в~DE~55 для EMV (Europay/Mastercard/Visa, чиповый платёжный стандарт; подробно~--- в~гл.~\ref{ch:emv}), делают сообщение самоописывающимся, и~каждое поле несёт собственный тег и~длину.
Bitmap проектировался в 1970-х~-- 80-х, когда каждый байт в~телекоммуникационном канале имел цену, и~сдвинуть бит для проверки присутствия поля было дешевле, чем пройти по тегам. К~моменту, когда форматы TLV с~сопоставимой производительностью стали широко доступны, выбор уже был закреплён стандартом.

\begin{figure}[tbp]
  \centering
  \includefiguresvg[width=.75\linewidth]{assets/figures/ch05-iso8583/bitmap-to-fields}
  \caption{От bitmap к~полям данных.}\label{fig:bitmap-to-fields}
\end{figure}

В~нашем примере основная bitmap выглядит так:

\begin{center}
  \texttt{72 3C 04 81 20 C0 80 00}
\end{center}

Переведём байты в~двоичную форму:

\begin{center}
  \small
  \texttt{72 = 0111\,0010} \quad
  \texttt{3C = 0011\,1100} \quad
  \texttt{04 = 0000\,0100} \quad
  \texttt{81 = 1000\,0001}

  \texttt{20 = 0010\,0000} \quad
  \texttt{C0 = 1100\,0000} \quad
  \texttt{80 = 1000\,0000} \quad
  \texttt{00 = 0000\,0000}
\end{center}

Первый бит равен \texttt{0}, так что дополнительной битовой карты нет и~всё сообщение укладывается в~первые 64 элемента данных. Остальные единицы указывают на DE~2, DE~3, DE~4, DE~7, DE~11, DE~12, DE~13, DE~14, DE~22, DE~25, DE~32, DE~35, DE~41, DE~42 и DE~49.
Это правдоподобный набор для обычной авторизации, в~котором есть идентификация карты, сумма, временной контекст, точка приёма и~способ ввода.

Bitmap отделяет структуру от содержания. Значения DE~4 и DE~41 ещё не прочитаны, но уже известно, что они есть и~идут в~порядке возрастания номеров.

Этот же разбор формализуется компактным декодером. Он принимает байты
сообщения, возвращает MTI и список номеров присутствующих DE; чтение
самих полей оставлено за рамками примера~--- для каждого DE нужны
его длина и формат из IFS партнёра.

\codefile{Python}{samples/ch05-iso8583-bitmap/bitmap.py}

\practicenote{%
  При разборе инцидента сначала смотрите на bitmap. Он быстро
  показывает, есть ли в~сообщении DE~55, Track 2, DE~38 и DE~39. Часто
  этого хватает, чтобы снять половину вопросов ещё до полного чтения дампа.

}
\section{Элементы данных -- пять групп вместо каталога}\label{sec:iso8583-data-elements}

Полный каталог элементов данных здесь не нужен. Важнее, какую работу поля выполняют внутри сообщения.
Для MTI~0100/MTI~0110 практически значимое раскладывается в~пять смысловых групп; полный каталог содержится в~сетевой спецификации и в~IFS партнёра~\cite{fis-iso8583-guide-2023}.

\subsection*{Идентификация и маршрутизация}

DE~2 (PAN) говорит, чья карта. DE~35 (Track~2 Data) дублирует ответ в~более богатой форме и~укладывает PAN, срок действия, сервисный код и дискреционные данные в~одну строку.
DE~32~-- идентификатор эквайера или банка-инициатора. DE~37 (RRN, Retrieval Reference Number~--- сквозной ссылочный номер операции) появляется в~ответе, в~запросах и~в~последующих сообщениях, когда операцию приходится искать в~отчётах, диспутах и операторских инструментах~\cite{fis-iso8583-guide-2023}.

Группа адресации отвечает, кому принадлежит карта, кто отправил запрос и~по какому идентификатору ту же операцию найдут в~другом контексте.

\subsection*{Деньги и~время}

DE~3 (Processing Code) кодирует тип операции: покупка, возврат, снятие наличных, запрос баланса.
DE~4 хранит сумму с~фиксированным масштабом, который определяется кодом валюты из DE~49 по таблице экспонентов ISO~4217 (для рубля и~евро~-- две минорные единицы, для иены~-- ноль).
DE~7 фиксирует сетевое время передачи сообщения. DE~11~-- STAN (Systems Trace Audit Number), шестизначный номер трассировки, по которому процессинг связывает запрос с~ответом. DE~12 и DE~13~-- локальное время и~дата терминала. DE~49~-- код валюты.

Группа нужна для сверки и расследований. Сумма без валюты ничего не значит, STAN без временного контекста не уникален, локальное время без сетевого~-- ненадёжный ориентир в~споре между терминалом, процессингом и эмитентом.

\subsection*{Контекст точки приёма}

DE~22 (POS Entry Mode) показывает, как получены данные карты: вручную, с~магнитной полосы, с~чипа, бесконтактно, в~интернет-операции. DE~25 (POS Condition Code) добавляет условия операции. DE~41 и DE~42~-- идентификаторы терминала и ТСП~-- привязывают транзакцию к~точке приёма.

Эта группа отвечает эмитенту, \emph{где} и \emph{как} прошла операция. Одни и~те же PAN и~сумма получают разную оценку, когда речь идёт о~чиповой покупке в~супермаркете или о~ручном вводе реквизитов на сайте.

\subsection*{Результат авторизации}

DE~39 (Response Code)~-- главное поле ответа, говорящее, одобрена операция или нет. DE~38 (идентификатор авторизации) появляется при одобрении и связывает авторизацию с~последующим клирингом.

Группа принципиально \emph{ответная}, в~запросе её нет, потому что решение ещё не принято. Чтение MTI~0110 начинается с DE~39, в~котором и~сидит смысл ответа.

\subsection*{Чиповые данные}

DE~55~-- контейнер TLV для данных EMV, в~котором лежат криптограммы, TVR (Terminal Verification Results, результаты проверок терминала), AID (Application Identifier, идентификатор платёжного приложения карты), ATC (Application Transaction Counter, счётчик операций приложения) и~другие теги, которыми терминал доказывает эмитенту, что операция пришла с~чипа или бесконтактного ядра. Подробный разбор тегов~--- в~гл.~\ref{ch:emv}~\cite{emvco-specs}.

В~каноническом примере DE~55 нет умышленно. Без чипового хвоста проще разделить грамматику ISO~8583 и~детали EMV. В~реальной чиповой авторизации именно DE~55 делает сообщение длиннее и~сложнее.

\section{Разбираем MTI~0100 и MTI~0110 до конца}\label{sec:iso8583-hex-dump}

MTI~0100 целиком:

\begin{center}
  \small
  \texttt{30 31 30 30}\\[2pt]
  \texttt{72 3C 04 81 20 C0 80 00}\\[2pt]
  \texttt{31 36 32 32 30 30 31 32 31 32 33 34 35 36 30 38 39 32}\\[2pt]
  \texttt{30 30 30 30 30 30}\\[2pt]
  \texttt{30 30 30 30 30 30 31 32 35 30 35 30}\\[2pt]
  \texttt{30 33 31 37 31 37 33 30 34 35}\\[2pt]
  \texttt{34 38 32 37 33 31}\\[2pt]
  \texttt{31 37 33 30 34 35}\\[2pt]
  \texttt{30 33 31 37}\\[2pt]
  \texttt{32 36 31 32}\\[2pt]
  \texttt{30 37 31}\\[2pt]
  \texttt{30 30}\\[2pt]
  \texttt{30 36 31 32 33 34 35 36}\\[2pt]
  \texttt{33 34 32 32 30 30 31 32 31 32 33 34 35 36 30 38 39 32}\\[2pt]
  \texttt{3D 32 36 31 32 32 30 31 31 32 33 34 35 36 37 38 39 30}\\[2pt]
  \texttt{54 45 52 4D 30 30 30 31}\\[2pt]
  \texttt{4D 45 52 43 48 41 4E 54 30 30 30 30 30 30 31}\\[2pt]
  \texttt{36 34 33}
\end{center}

\textbf{Шаг 1: MTI.} Первые четыре байта \texttt{30 31 30 30} в ASCII дают \texttt{0100}~-- запрос на авторизацию.

\textbf{Шаг 2: bitmap.} Следующие восемь байт \texttt{72 3C 04 81 20 C0 80 00} перечисляют поля: DE~2, DE~3, DE~4, DE~7, DE~11, DE~12, DE~13, DE~14, DE~22, DE~25, DE~32, DE~35, DE~41, DE~42, DE~49. Первый бит не установлен, и~вторичной битовой карты в~сообщении нет.

\textbf{Шаг 3: первое переменное поле.} DE~2 идёт как поле LLVAR (Length-Length-Variable~--- поле переменной длины с~двухсимвольным префиксом длины). Два байта длины после bitmap, \texttt{31 36}, в ASCII дают \texttt{16}; дальше 16 цифр PAN~-- \texttt{2200121234560892}.

Дамп разбирается по логике одного конкретного профиля. Точные коды DE~22, DE~25, список обязательных полей и~пара сообщений отмены зависят от спецификации интерфейса. Глава даёт модель чтения; универсальным справочником всех диалектов она не служит~\cite{fis-iso8583-guide-2023}.

\textbf{Шаг 4: остальные поля~-- в~порядке возрастания номеров.}
Сразу после PAN идут DE~3=\texttt{000000} и DE~4=\texttt{000000125050}. Код операции означает стандартную покупку, сумма~-- 1\,250 рублей 50~копеек при DE~49=\texttt{643}.
Дальше DE~7=\texttt{0317173045}, DE~11=\texttt{482731}, DE~12=\texttt{173045}, DE~13=\texttt{0317}. Это временные метки операции, в~которых есть сетевое время, STAN для связи запроса с~ответом, локальные дата и~время терминала.

DE~14=\texttt{2612}~-- срок действия карты. DE~22=\texttt{071} в~этом профиле описывает бесконтактный чиповый сценарий с~терминалом, способным принять PIN (третий разряд кодирует PIN entry capability терминала; факт верификации PIN в~этой операции через него не передаётся). DE~25=\texttt{00} означает обычные условия операции.
За ними идут поля маршрутизации и точки приёма. DE~32=\texttt{123456} указывает на эквайера или банка-инициатора; DE~35 несёт Track~2 с PAN, сроком действия и сервисным кодом; DE~41=\texttt{TERM0001} и DE~42=\texttt{MERCHANT0000001} идентифицируют терминал и~ТСП. Последним приходит DE~49=\texttt{643}~-- код рубля.

MTI~0100 прочитан целиком: \enquote{эквайер просит одобрить бесконтактную покупку на сумму 1\,250 рублей 50 копеек по карте с PAN \texttt{2200121234560892} в~терминале \texttt{TERM0001} у~ТСП \texttt{MERCHANT0000001}}.

Эмитент отвечает MTI~0110. Ответ оставляет только тот контекст, который достаточен для сопоставления с~запросом, и добавляет результат.
MTI меняется на MTI~0110, bitmap становится короче, например \texttt{32 38 00 00 0E C0 80 00}, то есть DE~3, DE~4, DE~7, DE~11, DE~12, DE~13, DE~37, DE~38, DE~39, DE~41, DE~42, DE~49. PAN и~Track~2 уже не нужны; к~эху сумм, времени и~идентификаторов точки приёма добавляются поля результата.
DE~37 несёт RRN~-- сквозной ссылочный номер для сопоставления в~клиринге. DE~38 приносит код авторизации, скажем \texttt{A1B2C3}. DE~39 фиксирует решение эмитента: \texttt{00}~-- одобрение, \texttt{51}~-- недостаток средств. Это типичный формат ответа в~одном из публично документированных профилей~\cite{fis-iso8583-guide-2023}.
Сообщение отмены~-- MTI~0400 или MTI~0420 в~зависимости от IFS~-- должно нести ссылочные данные на исходную авторизацию: тот же STAN, тот же временной контекст; в~ряде профилей~-- специальные original data elements~\cite{fis-iso8583-guide-2023}.

\importantnote{%
  В~примере выше все числовые поля закодированы в ASCII, и~каждая
  цифра занимает один байт. В~реальных IFS цифры могут кодироваться
  упакованным числовым форматом (BCD, Binary Coded Decimal), который
  кладёт две цифры в~один байт и~вдвое сжимает числовые поля.
  ASCII удобнее при отладке в~логах, а~BCD типичен для продуктовых (\textit{production}) систем,
  где каждый лишний байт увеличивал стоимость телеком-канала и~время
  обработки на железе 1980-х. Прежде чем читать дамп вручную, нужно
  узнать кодировку из IFS, иначе ошибка на этом этапе делает бессмысленным
  весь дальнейший разбор~\cite{fis-iso8583-guide-2023}.

}
\section{Где заканчивается ISO~8583 и~зачем нужна IFS}\label{sec:iso8583-implementations}

ISO~8583 не описывает сетевой протокол целиком. Стандарт определяет \emph{грамматику сообщения}; в~официальной аннотации ISO прямо оговорено, что метод расчёта в~предмет стандарта не входит.
Заметная часть реальной интеграции лежит вне этой грамматики~\cite{iso-8583-2023,fis-iso8583-guide-2023}.

\textbf{Транспорт и фрейминг.} Способ доставки~-- по TCP, через VPN, с~каким заголовком и префиксом длины~-- стандартом не предписан. Двух- или четырёхбайтовый префикс длины, отдельный application header, выбор между ASCII и EBCDIC (Extended Binary Coded Decimal Interchange Code~--- 8-битная кодировка IBM, исторически используемая на мейнфреймах процессингов)~-- всё это соглашение конкретного процессинга.
Транспортный нейтралитет был проектным выбором: стандарт 1987 года должен был одинаково работать на X.25, SNA и~любом будущем протоколе. Цена этого решения в~том, что каждый IFS фиксирует транспорт сам и~при смене партнёра ничего из этого не переносится.

\textbf{Криптография и~ключи.} PIN block (зашифрованный блок с~PIN-кодом фиксированной длины; форматы~--- в~§\,\ref{sec:crypto-pin-block}) едет в~DE~52, MAC (Message Authentication Code, код аутентификации сообщения)~--- в~DE~64 или DE~128. Алгоритмы, формат PIN block, схема ключей и~порядок их обмена задаются сетью или спецификацией конкретного коммутатора; сам ISO~8583 их не описывает.
\highlight{Само наличие DE~64 или DE~128 не объясняет MAC}: для проверки
нужны ключи, алгоритм, padding, состав строки и транспортная
спецификация~--- обвязка вокруг ISO~8583, лежащая за пределами его
базовой грамматики.

\textbf{Состояние сессии.} Sign-on, echo test, cutover, управление ключами и повторная отправка опираются на MTI~0800/MTI~0810 и~на то, как конкретная система понимает таймауты, идемпотентность (гарантию, что повторная отправка того же сообщения не приведёт к~двойному списанию или двойной отмене) и повтор.

\textbf{Правила обработки.} Что считать дубликатом, когда отправлять сообщение отмены, как долго ждать ответ, какие поля обязательны для чиповой или интернет-операции~--- всё это прикладные правила интерфейса, лежащие за рамками синтаксиса сообщения.

Прочитанного канонического запроса и~ответа для подключения нового процессинга мало. Стандарт задаёт общий язык, а~интеграция всегда идёт на одном из его диалектов.

IFS (Interface Specification)~-- спецификация интерфейса конкретной сети, процессинга или банка-партнёра. Она фиксирует обязательные поля, кодировки, правила повторов и транспортную обвязку канала~\cite{fis-iso8583-guide-2023}.

Диалекты расходятся сразу по нескольким направлениям. В~одной системе всё кодируется в ASCII, в~другой рядом живут ASCII, BCD и~бинарные поля. Одному партнёру достаточно DE~35, другому нужен ещё DE~2, третьему~-- набор частных полей под конкретный сценарий. Где-то дополнительные данные раскрываются через DE~48, DE~60, DE~61, DE~63 или DE~124, где-то собираются в~собственные подполя. Аналогично расходятся коды ответа, жизненный цикл сообщений отмены, cutover, транспортный заголовок, sign-on и~ключевой обмен.

Одной книги по ISO~8583 для интеграции мало. Стандарт описывает грамматику, а~подключение всегда упирается в IFS конкретной сети или процессинга. Только IFS говорит, какие поля обязательны для вашего канала, как кодируется сумма, где искать частные подполя и~каким должен быть ответ на ваш MTI~0800.

Детали диалектов Visa, Mastercard и НСПК здесь не разбираются. Полный справочник элементов данных и~дампы MTI~0200, MTI~0420 и~сетевых служебных сообщений живут в~сетевой спецификации и в IFS; книга оставляет за собой устойчивую модель чтения.
